\documentclass[letterpaper,10pt,english]{article}
\usepackage{%
amsfonts,%
amsmath,%
amssymb,%
amsthm,%
babel,%
bbm,%
%biblatex,%
caption,%
centernot,%
color,%
enumerate,%
%enumitem,%
epsfig,%
epstopdf,%
etex,%
fancybox,%
framed,%
fullpage,%
%geometry,%
graphicx,%
hyperref,%
latexsym,%
mathptmx,%
mathtools,%
multicol,%
paralist,%
pgf,%
pgfplots,%
pgfplotstable,%
pgfpages,%
proof,%
psfrag,%
%subfigure,%
tcolorbox,%
tfrupee,%
tikz,%
times,%
%ulem,%
url,%
xcolor,%
mathpazo,%
}

\definecolor{shadecolor}{gray}{.95}%{rgb}{1,0,0}
\usepackage[margin=1in,top=0.75in]{geometry}
\usepackage[mathscr]{eucal}
\usepgflibrary{shapes}
\usepgfplotslibrary{fillbetween}
\usetikzlibrary{%
arrows,%
backgrounds,%
chains,%
decorations.pathmorphing,% /pgf/decoration/random steps | erste Graphik
decorations.text,% 
matrix,%
positioning,% wg. " of "
fit,%
patterns,%
petri,%
plotmarks,%
scopes,%
shadows,%
shapes.misc,% wg. rounded rectangle
shapes.arrows,%
shapes.callouts,%
shapes%
}

%\pgfplotsset{compat=newest} %<------ Here
\pgfplotsset{compat=1.11} %<------ Or use this one

\theoremstyle{plain}
\newtheorem{thm}{Theorem}[section]
\newtheorem{lem}[thm]{Lemma}
\newtheorem{prop}[thm]{Proposition}
\newtheorem{cor}[thm]{Corollary}
\newtheorem{clm}[thm]{Claim}

\theoremstyle{definition}
\newtheorem{defn}[thm]{Definition}
\newtheorem{conj}[thm]{Conjecture}
\newtheorem{exmp}[thm]{Example}
\newtheorem{axiom}[thm]{Axiom}
\newtheorem{assum}[thm]{Assumption}
\newtheorem{exerc}[thm]{Exercise}
\newtheorem*{defin}{Definition}

\theoremstyle{remark}
\newtheorem{rem}{Remark}
\newtheorem{note}{Note}

\newcommand{\iid}{\emph{i.i.d.}\ }
\newcommand{\Cov}{\operatorname{Cov}}
%\newcommand{\det}{\operatorname{det}}
%\newcommand{\Real}{\mathbb{R}}
\newcommand{\tr}{\operatorname{tr}}
\newcommand{\Var}{\operatorname{Var}}
\newcommand{\cov}{\operatorname{cov}}

%\renewcommand{\proof}[1]{\begin{proof}#1\end{proof}}
\newcommand{\EQ}[1]{\begin{equation*}#1\end{equation*}}
\newcommand{\EQN}[1]{\begin{equation}#1\end{equation}}
\newcommand{\eq}[1]{\begin{align*}#1\end{align*}}
\newcommand{\meq}[2]{\begin{xalignat*}{#1}#2\end{xalignat*}}
\newcommand{\norm}[1]{\left\lVert#1\right\rVert}
\newcommand{\inner}[1]{\left\langle #1\right\rangle}
\newcommand{\abs}[1]{\left\lvert#1\right\rvert}
\newcommand{\expect}[1]{\mathbb{E}\left[{#1}\right]}
\newcommand{\prob}[1]{\mathbb{P}\left[{#1}\right]}
\newcommand{\given}{\; \big\vert \;} 
\newcommand{\set}[1]{\left\{#1\right\}} 
\newcommand{\indicator}[1]{1_{\set{#1}}} 
\newcommand{\Ind}[1]{\mathbbm{1}_{#1}} 
\newcommand{\SetIn}[1]{\mathbbm{1}_{\set{#1}}} 
\newcommand{\red}[1]{\textcolor{red}{#1}} 
\newcommand{\floor}[1]{\left\lfloor#1\right\rfloor}
\newcommand{\ceil}[1]{\left\lceil#1\right\rceil}


\newcommand{\C}{\mathbb{C}}
\newcommand{\D}{\mathbb{D}}
\newcommand{\E}{\mathbb{E}}
\newcommand{\F}{\mathbb{F}}
\newcommand{\N}{\mathbb{N}}
\renewcommand{\P}{\mathbb{P}}
\newcommand{\Q}{\mathbb{Q}}
\newcommand{\R}{\mathbb{R}}
\newcommand{\Z}{\mathbb{Z}}

\newcommand{\bA}{\mathbf{A}}
\newcommand{\bI}{\mathbf{I}}
\newcommand{\bU}{\mathbf{1}}
\newcommand{\bx}{\mathbf{x}}
\newcommand{\bX}{\mathbf{X}}
\newcommand{\bZ}{\mathbf{Z}}


\newcommand{\cA}{\mathcal{A}}
\newcommand{\cB}{\mathcal{B}}
\newcommand{\cC}{\mathcal{C}}
\newcommand{\cF}{\mathcal{F}}
\newcommand{\cG}{\mathcal{G}}
\newcommand{\cM}{\mathcal{M}}
\newcommand{\cN}{\mathcal{N}}
\newcommand{\cP}{\mathcal{P}}
\newcommand{\cT}{\mathcal{T}}
\newcommand{\cX}{\mathcal{X}}
\newcommand{\cY}{\mathcal{Y}}

\newcommand{\sA}{\mathscr{A}}
\newcommand{\sB}{\mathscr{B}}
\newcommand{\sC}{\mathscr{C}}
\newcommand{\sE}{\mathscr{E}}
\newcommand{\sF}{\mathscr{F}}
\newcommand{\sG}{\mathscr{G}}
\newcommand{\sH}{\mathscr{H}}
\newcommand{\sL}{\mathscr{L}}
\newcommand{\sS}{\mathscr{S}}
\newcommand{\sT}{\mathscr{T}}
\newcommand{\sX}{\mathscr{X}}
\newcommand{\sY}{\mathscr{Y}}
\newcommand{\sZ}{\mathscr{Z}}

\newcommand{\tS}{\tilde{S}}
% Debug
\newcommand{\todo}[1]{\begin{color}{blue}{{\bf~[TODO:~#1]}}\end{color}}

% a few handy macros

\renewcommand{\le}{\leqslant}
\renewcommand{\ge}{\geqslant}
\newcommand\matlab{{\sc matlab}}
\newcommand{\goto}{\rightarrow}
\newcommand{\bigo}{{\mathcal O}}
%\newcommand{\half}{\frac{1}{2}}
%\newcommand\implies{\quad\Longrightarrow\quad}
\newcommand\reals{{{\rm l} \kern -.15em {\rm R} }}
\newcommand\complex{{\raisebox{.043ex}{\rule{0.07em}{1.56ex}} \hskip -.35em {\rm C}}}


% macros for matrices/vectors:

% matrix environment for vectors or matrices where elements are centered
\newenvironment{mat}{\left[\begin{array}{ccccccccccccccc}}{\end{array}\right]}
\newcommand\bcm{\begin{mat}}
\newcommand\ecm{\end{mat}}

% matrix environment for vectors or matrices where elements are right justifvied
\newenvironment{rmat}{\left[\begin{array}{rrrrrrrrrrrrr}}{\end{array}\right]}
\newcommand\brm{\begin{rmat}}
\newcommand\erm{\end{rmat}}

% for left brace and a set of choices
%\newenvironment{choices}{\left\{ \begin{array}{ll}}{\end{array}\right.}
\newcommand\when{&\text{if~}}
\newcommand\otherwise{&\text{otherwise}}
% sample usage:
%\delta_{ij} = \begin{choices} 1 \when i=j, \\ 0 \otherwise \end{choices}


% for labeling and referencing equations:
\newcommand{\eql}{\begin{equation}\label}
\newcommand{\eqn}[1]{(\ref{#1})}
% can then do
%\eql{eqnlabel}
%...
%\end{equation}
% and refer to it as equation \eqn{eqnlabel}.
% some useful macros for finite difference methods:
\newcommand\unp{U^{n+1}}
\newcommand\unm{U^{n-1}}
% for chemical reactions:
\newcommand{\react}[1]{\stackrel{K_{#1}}{\rightarrow}}
\newcommand{\reactb}[2]{\stackrel{K_{#1}}{~\stackrel{\rightleftharpoons}
 {\scriptstyle K_{#2}}}~}
\makeatletter
\def\th@plain{%
\thm@notefont{}% same as heading font
\itshape % body font
}
\def\th@definition{%
\thm@notefont{}% same as heading font
\normalfont % body font
}
\makeatother
\date{}



\title{ Lecture-21: Recurrent and transient states}
\author{}

\begin{document}
\maketitle


\section{Recurrence and Transience}
We will consider a random sequence $X:\Omega\to\sX^{\Z_+}$ with initial state $X_0 = x\in\sX$ and the $k$th hitting times to state $y$ for all $k \in \N$ denoted by $\tau_k \triangleq \tau_X^{\set{y},k}$ inductively defined as $\tau_k \triangleq \inf\set{n > \tau_{k-1}: X_n = y}$ where $\tau_0 \triangleq 0$. 
We define the inter-return time sequence $H:\Omega\to\N^\N$ as $H_k \triangleq H_X^{\set{y},k} = \tau_k - \tau_{k-1}$ for all $k \in \N$. 
%\begin{defn} 
%For a process $X:\Omega\to\sX^{\R_+}$ event $A$ and random variable $Y:\Omega\to\R$, 
%we denote the conditional probability and conditional distribution, given the initial state $\set{X_0 = x}$ by $P_x(A) \triangleq P(A\mid\set{X_0 = x})$ and $\E_xY = \E[Y\mid\set{X_0 = x}]$ respectively.
%\end{defn}
\begin{defn} 
For a random sequence $X: \Omega \to \sX^{\Z_+}$ with initial state $X_0 = x$, 
\begin{compactenum}[(i)]
\item the \textbf{probability of hitting state $y$ eventually} is denoted by 
$f_{xy} \triangleq P_x\set{\tau_1 < \infty}$, 
and 
\item the \textbf{probability of first visit to state $y$ at time $n \in\N$} is denoted by 
$f_{xy}^{(n)} 
\triangleq P_x\set{\tau_1 = n}. %,\quad n \in \N.
$
\end{compactenum}
\end{defn}
%\begin{defn}
%The probability that the Markov chain $X$ hits state $y$ eventually, starting from initial state $x$ is 
%\EQ{
%%f_{xy} \triangleq P(\set{H_y < \infty} \given \set{X_0 = x}) = P(\cup_{n \in \N}\set{H_y = n} \given \set{X_0 = x}) = \sum_{n \in \N}P(\set{H_y = n}\given \set{X_0 = x}) = \sum_{n \in \N}f_{xy}^{(n)}.
%f_{xy} 
%\triangleq P_x\set{H_1 < \infty} 
%= P_x(\cup_{n \in \N}\set{H_1 = n}) 
%= \sum_{n \in \N}P_x\set{H_1 = n} 
%= \sum_{n \in \N}f_{xy}^{(n)}.
%}
%\end{defn}
\begin{rem}
We can write the finiteness of hitting time $\tau_1$ as the disjoint union $\set{\tau_1 < \infty} = \cup_{n \in \N}\set{\tau_1=n}$. 
Therefore, $f_{xy} = \sum_{n\in\N}f_{xy}^{(n)}$. 
\end{rem}
\begin{rem}
If $f_{xy} = P_x\set{\tau_1 < \infty} = 1$ for all initial states $x \in \sX$, 
then $\tau_1$ is almost surely finite and hence a stopping time.
\end{rem}
\begin{defn} 
From the initial state $x$, the distribution 
\begin{compactenum}[(i)]
\item for the first hitting time to state $y$ is called the \textbf{first passage time distribution} and denoted by $((f_{xy}^{(n)}: n \in \N), 1 - f_{xy})$, and 
\item for the first return time to state $x$ is called the \textbf{first recurrence time distribution} and denoted by $((f_{xx}^{(n)}: n \in \N),1 - f_{xx})$. 
\end{compactenum}
\end{defn}
\begin{defn}
A state $y\in\sX$ is called \textbf{recurrent} if $f_{yy} = 1$, and is called \textbf{transient} if $f_{yy} < 1$. 
\end{defn}
\begin{defn}
For any state $y \in \sX$, the \textbf{mean recurrence time} is denoted by $\mu_{yy} \triangleq \E_y\tau_1$. 
\end{defn}
\begin{rem}
The mean recurrence time for any transient state is infinite. 
For any recurrent state $y \in \sX$, we write $\tau_1 = \tau_1\SetIn{\tau_1 < \infty} = \sum_{n \in \N}n\SetIn{\tau_1=n}$ almost surely, 
and the mean recurrence time is given by $\mu_{yy} = \sum_{n \in \N}nf^{(n)}_{yy}$.   
\end{rem}
\begin{defn} 
For a recurrent state $y \in \sX$, 
\begin{compactenum}[(i)]
\item if the mean recurrence time is finite, then the state $y$ is called \textbf{positive recurrent}, and  
\item if the mean recurrence time is infinite, then the state $y$ is called \textbf{null recurrent}.
\end{compactenum}
\end{defn}

%\begin{defn} 
%Let $\tau_y^{(0)} = 0$ and define $\tau_k$ to be the \textbf{$k$th hitting time of state $y$}, defined as 
%\EQ{
%\tau_k \triangleq \inf\set{n > \tau_X^{\set{y},k-1}: X_n = y}, ~ k \in \N.
%}
%We can also define the \textbf{$k$th excursion time} as $H_k \triangleq \tau_k-\tau_X^{\set{y},k-1}$, for all $k \in \N$. 
%\end{defn}


\begin{prop} 
For a homogeneous discrete Markov chain $X: \Omega \to \sX^{\Z_+}$, 
%The total number of visits to a state $y \in \sX$ %after starting from initial state $x$ 
%is denoted by $N_y = \sum_{n \in \N}\indicator{X_n = y}$.
%Then, for each $m \in \Z_+$, 
we have 
\EQ{
P_x\set{N_y(\infty) = k} =\begin{cases}
1 - f_{xy}, & k = 0,\\
f_{xy}f_{yy}^{k-1}(1-f_{yy}), & k \in \N.
\end{cases}
}
\end{prop}
\begin{proof} 
We can write the event of zero visits to state $y$ as $\set{N_y(\infty) = 0} = \set{\tau_1 = \infty}$. 
Further, we can write the event of $m$ visits to state $y$ as 
\EQ{
\set{N_y(\infty) = k} 
= \set{\tau_k < \infty}\cap\set{\tau_{k+1}=\infty}
= \cap_{j=1}^k\set{H_j< \infty}\cap\set{H_{k+1} = \infty},\quad k \in \N.
}
Recall that $H:\Omega\to\N^\N$ is an independent random sequence with subsequence $(H_k: k \ge 2)$ identically distributed, 
with $P_x\set{H_j = n} = P_y\set{\tau_1 =n}$ for all $j \ge 2$.  
Therefore, we get 
\EQ{
P_x\set{N_y(\infty)=k} 
= P_x\set{H_1 < \infty}\prod_{j=2}^kP_x\set{H_j < \infty}P_x\set{H_{k+1}=\infty}
= f_{xy}f_{yy}^{k-1}(1-f_{yy}).
}
%Conditioned on $X_0 = x$, the first passage time $H_y$ to state $y$ being finite is a Bernoulli random variable with probability $f_{xy}$. 
%The time of the $m$th return to the state $y$ is a recurrence time for each $m \in \Z_+$.  
%From strong Markov property, each return to state $y$ is independent of the past. 
%Hence, each return to state $y$ in a finite time is an \textit{iid} Bernoulli random variable with probability $f_{yy}$. 
%It follows that the number of recurrences to state $y$ is the time for first failure to return. 
%Conditioned on initial state being $X_0 = y$, the distribution of $N_y$ is geometric random variable with failure probability $1-f_{yy}$. 
%%Hence the probabilities will multiply. 
%%Hence starting from state $y$ we visit the state $y$ with probability $f_{yy}$. 
%%Moreover, we visit the state $y$ exactly $m$ number of times (and no more) with probability $f_{yy}^m(1-f_{yy}).$\\
%%Conditioned on $X_0 = x$, the stopping time of first visit to $y$ is a Bernoulli random variable with probability $f_{xy}$. 
%%Hence, the second result follows. 
\end{proof}
%\begin{proof}
%For $m=0$, we can write the following equality from the definition, 
%\EQ{
%P_x\set{N_y = 0} = P_x\set{H_1 = \infty} = 1- f_{xy}.
%} 
%For $m \in \N$, we first define events $E_k \triangleq \set{\tau_k < \infty}$ for all $k \in \N$. 
%%we consider $P_x\set{N_y > m}$. 
%%Let $\tau_y^{(0)} = 0$ and define $\tau_k$ to be the $k$th hitting time of state $y$, defined as 
%%%\EQ{
%%$\tau_k \triangleq \inf\set{n > \tau_X^{\set{y},k-1}: X_n = y}. $
%%%}
%%Then, we define the excursion times as $H_k \triangleq \tau_k-\tau_X^{\set{y},k-1}$, 
%%and define events $E_k \triangleq \set{\tau_k < \infty}$ for all $k \in \N$. 
%Then, 
%\EQ{
%P_x\set{N_y = m}  = P_x(\set{\tau_y^{(m)} < \infty}\cap\set{\tau_y^{m+1} = \infty}) = P_x(E_m\cap E_{m+1}^c).
%%= P_x(\set{\tau_y^{(m)} < \infty}\cap\set{H_y^{m+1} = \infty}).
%}
%%From the strong Markov property of homogeneous DTMC, we can write 
%%\eq{
%%P_x(\set{\tau_y^{(m)} < \infty}\cap\set{H_y^{m+1} = \infty}) &= 
%%P(\set{H_y^{m+1} = \infty} \given\set{X_0 = x, S_{y}^{(m)} < \infty, X_{\tau_y^{(m)}} = y})P_x(\set{\tau_y^{(m)} < \infty})\\
%%&= P_y(\set{H_y^{m+1} = \infty}P_x(\set{\tau_y^{(m)} < \infty}).
%%}
%We observe that $(E_k: k \in \N)$ is a decreasing sequence of events, and hence $E_m = \cap_{k=1}^mE_k$.  
%Together with this observation and from the definition of conditional probability, 
%we can write 
%\EQ{
%P_x\set{N_y = m} = P_x(E_1\cap E_2 \cap \dots \cap E_m \cap E_{m+1}^c) = P_x(E_1)\left(\prod_{k=2}^{m}P(E_k\given E_1 \cap \dots \cap E_{k-1}) \right)P(E_{m+1}^c\given E_1\cap \dots \cap E_m). 
%}
%From the definition, we get $P_x(E_1) = P_x\set{H_y < \infty} = f_{xy}$. 
%We focus on the conditional probability of the following event for $k \in \set{2, \dots, m}$, 
%which equals 
%\EQ{
%P_x(E_k \given E_1\cap \dots\cap E_{k-1}) 
%= P_x(E_k \given E_{k-1}) = P(\set{H_k< \infty}\given\set{ X_{\tau_X^{\set{y},k-1}} = y, \tau_X^{\set{y},k-1}< \infty, X_0 = x}) 
%= P_y\set{H_k< \infty} = f_{yy}. 
%}
%Equality in the second line follows from the strong Markov property and the definition of $f_{yy}$. 
%The result follows from the aggregation of the above equalities. 
%\end{proof}
%\begin{rem}
%For a time homogeneous Markov chain $X: \Omega\to\sX^{\Z_+}$ with initial state $x$, 
%\begin{compactenum}[(i)]
%\item $P_x\set{N_y < \infty} = \SetIn{f_{yy} < 1} + (1-f_{xy})\SetIn{f_{yy} = 1}$,
%\item 
%\end{compactenum}
%\end{rem}
\begin{cor}
\label{cor:01law}
For a homogeneous Markov chain $X$, we have 
%\eq{
$P_x\set{N_y(\infty) < \infty} = \SetIn{f_{yy} < 1} + (1-f_{xy})\SetIn{f_{yy} = 1}$. 
%}
\end{cor}
\begin{proof}
We can write the event $\set{N_y(\infty) < \infty}$ as disjoint union of events $\set{N_y(\infty) = k}$, to get the result. 
%\EQ{
%P_x\set{N_y < \infty} = \sum_{n \in\Z_+}P_x\set{N_y = n} = \SetIn{f_{yy} < 1} + (1-f_{xy})\SetIn{f_{yy} = 1}.
%}
\end{proof}
\begin{rem} 
For a time homogeneous Markov chain $X: \Omega \to \sX^{\Z_+}$, we have 
\begin{compactenum}[(i)]
\item $P_x\set{N_y(\infty) = \infty} = f_{xy}\SetIn{f_{yy}=1}$, and 
\item $P_y\set{N_y(\infty) = \infty} = \SetIn{f_{yy}=1}$.
\end{compactenum} 
\end{rem}
\begin{cor} 
The mean number of visits to state $y$, starting from a state $x$ is 
$
\E_xN_y(\infty) = \frac{f_{xy}}{1-f_{yy}}\SetIn{f_{yy} < 1} + \infty\SetIn{f_{xy} > 0, f_{yy} = 1}.
$
\end{cor}
\begin{rem} 
For any state $y \in \sX$, we have $\E_yN_y(\infty) = \frac{f_{yy}}{1-f_{yy}}\SetIn{f_{yy}<1} + \infty\SetIn{f_{yy}=1}$. 
That is, the mean number of visits to initial state $y$ is finite iff the state $y$ is transient.  
\end{rem}
%\begin{proof} 
%When $i=j$, it follows from the previous Proposition that $N_y$ is geometric random variable with success probability $1-f_{yy}$. 
%\end{proof}

\begin{shaded}
\begin{rem}
In particular, this corollary implies the following consequences. 
\begin{enumerate}[i\_]
	\item A transient state is visited a finite amount of times almost surely. 
	This follows from Corollary~\ref{cor:01law}, since $P_x\set{N_y(\infty) < \infty} = 1$ for all transient states $y \in \sX$ and any initial state $x \in \sX$. 
	\item A recurrent state is visited infinitely often almost surely. 
	This also follows from Corollary~\ref{cor:01law}, since $P_y\set{N_y(\infty) < \infty} = 0$ for all recurrent states $y \in \sX$.  
	\item In a finite state Markov chain, not all states may be transient. 
	\begin{proof}
	To see this, we assume that for a finite state space $\sX$, all states $y \in \sX$ are transient. 
	Then, we know that $N_y(\infty)$ is finite almost surely for all states $y \in \sX$. 
	It follows that, for any initial state $x \in \sX$
	\EQ{
	0 \le P_x\set{\sum_{y \in \sX}N_y(\infty) = \infty} = P_x(\cup_{y \in \sX}\set{N_y(\infty) = \infty}) \le \sum_{y \in \sX}P_x\set{N_y(\infty) = \infty} = 0. 
	}
It follows that $\sum_{y \in \sX}N_y(\infty)$ is also finite almost surely for all states $y \in \sX$ for finite state space $\sX$. 
However, we know that $\sum_{y \in \sX}N_y(\infty) = \sum_{k \in \N}\sum_{y \in \sX}\indicator{X_k = y} = \infty$. 
This leads to a contradiction. 
\end{proof}
\end{enumerate} 
\end{rem}
\end{shaded}
\begin{prop}
For a homogeneous DTMC $X: \Omega \to \sX^{\Z_+}$, a state $y\in\sX$ is recurrent iff
%\begin{align*}
$\sum_{k\in\N} p_{yy}^{(k)} = \infty$, and transient iff $\sum_{k\in\N} p_{yy}^{(k)} < \infty$.
%\end{align*}
\end{prop}
\begin{proof}
%For any state $x \in \sX$, we can write $p_{xx}^{(k)} = P_x\set{X_k=x} = \E_x \SetIn{X_k=x}$. 
%Using monotone convergence theorem to exchange expectation and summation, we obtain  
%$
%\sum_{k\in \N}p_{xx}^{(k)}=\E_x \sum_{k\in\N}\SetIn{X_k=x} = \E_x N_x.
%$
Recall that if the mean recurrence time to a state $y\in\sX$ is $\E_yN_y(\infty) = \sum_{k\in\N}p_{yy}^{(k)}$ finite then the state is transient and infinite if the state is recurrent.
\end{proof}
\begin{cor} 
For a transient state $y \in \sX$, the following limits hold 
$\lim_{n \to \infty}p_{xy}^{(n)} = 0$, and $\lim_{n \to \infty}\frac{\sum_{k = 1}^np_{xy}^{(k)}}{n} = 0$.
\end{cor}
\begin{proof} 
For a transient state $y \in \sX$ and any state $x \in \sX$, we have $\E_xN_y(\infty)  = \sum_{n \in \N}p_{xy}^{(n)} < \infty$. 
Since the series sum is finite, it implies that the limiting terms in the sequence $\lim_{n \to \infty}p_{xy}^{(n)} = 0$. 
Further, we can write $\sum_{k=1}^np_{xy}^{(k)} \le \E_xN_y(\infty) \le M$ for some $M \in \N$ and hence $\lim_{n \to \infty}\frac{\sum_{k = 1}^np_{xy}^{(k)}}{n} = 0$.
\end{proof}
%\begin{cor}
%For a transient state $y \in \sX$, we have $P(N_y < \infty) = 1$. 
%\end{cor}
%\begin{proof}
%\end{proof}
\begin{lem} 
\label{lem:StoppingTime}
For any state $y \in \sX$, let $H:\Omega\to\N^\N$ be the sequence of almost surely finite inter-visit times to state $y$, 
and $N_y(n) = \sum_{k=1}^n\indicator{X_k = y}$ be the number of visits to state $y$ in $n$ times. 
Then, $N_y(n) +1$ is a finite mean stopping time with respect to the sequence $H$. 
\end{lem}
\begin{proof} 
We first observe that $N_y(n) + 1 \le n+1$ and hence has a finite mean for each $n \in \N$. 
Further, we observe that $\set{N_y(n) + 1 = k}$ can be completely determined by observing $H_1, \dots, H_k$, since 
%To see this, we notice that 
\EQ{
\set{N_y(n) + 1 = k} 
= \set{\sum_{j=1}^{k-1}H_j \le n <  \sum_{j=1}^kH_j} 
\in \sigma(H_1, \dots, H_k). 
}
\end{proof}
%We define $N_y(n) \triangleq \sum_{k=1}^n\indicator{X_k = y}$ to be the number of visits to state $y$ in $n$ steps of the Markov process $X$. `
%Then, $\E_xN_y(n)  =  \sum_{k=1}^np_{xy}^{(k)}$.  
\begin{thm} 
Let $x, y \in \sX$ be such that $f_{xy} = 1$ and $y$ is recurrent.  
Then, $\lim_{n \to \infty}\frac{\sum_{k=1}^np_{xy}^{(k)}}{n} = \frac{1}{\mu_{yy}}$. 
%Furthermore, if $y$ is aperiodic, then $\lim_{n \to \infty}P_{xy}^{(n)} = \frac{1}{\mu_{yy}}$. 
\end{thm}
\begin{proof} 
Let $y \in \sX$ be recurrent. 
The proof consists of three parts. 
In the first two parts, we will show that starting from the state $y$, we have the limiting empirical average of mean number of visits to state $y$ is $\lim_{n \to \infty}\frac{1}{n}\E_yN_y(n) = \frac{1}{\mu_{yy}}$. 
In the third part, we will show that for any starting state $x \in \sX$ such that $f_{xy}=1$, we have  the limiting empirical average of mean number of visits to state $y$ is $\lim_{n \to \infty}\frac{1}{n}\E_xN_y(n) = \frac{1}{\mu_{yy}}$. 
\begin{itemize}
\item[Lower bound:]
%> \sum_{\ell=1}^{N_y(n)}H_1$. 
We observe that $N_y(n)+1$ is a stopping time with respect to inter-visit times $H$ from Lemma~\ref{lem:StoppingTime}. 
Further, we have $\sum_{j=1}^{N_y(n)+1}H_j > n$. 
Applying Wald's Lemma to the random sum $\sum_{j=1}^{N_y(n)+1}H_j$ , we get $\E_y(N_y(n)+1)\mu_{yy} > n$. 
Taking limits, we obtain $\lim\inf_{n \in \N}\frac{\sum_{k=1}^np_{yy}^{(k)}}{n} \ge \frac{1}{\mu_{yy}}$. 

\item[Upper bound:] 
Let $X_0 = y$ and consider a fixed positive integer $M\in \N$. 
Then $H$ is \iid and we define truncated recurrence times $\bar{H}:\Omega\to[M]^\N$ for all $j \in \N$ as  $\bar{H}_j \triangleq M \wedge H_j$. 
It follows that the sequence $\bar{H}$ is \iid and $\bar{H}_j \le H_j$ for all $j \in \N$. 
We define the mean of the truncated recurrence times as $\bar{\mu}_{yy} \triangleq \E_y\bar{H}_1$. 
From the monotonicity of truncation, we get $\bar{\mu}_{yy} \le \mu_{yy}$.  

We define the random variable $\bar{\tau}_k \triangleq \sum_{j=1}^k\bar{H}_j$ for all $k \in \N$, 
and $\bar{\tau}_k \le \tau_k$ for all $k \in \N$. 
We can define the associated counting process that counts the number of truncated recurrences in first $n$ steps as 
$\bar{N}_y(n) \triangleq \sum_{k\in \N}\SetIn{\bar{\tau}_k \le n}$ for all $n \in \N$. 
We conclude that $\bar{N}_y(n)+1$ is a stopping time with respect to \iid process $\bar{H}$, 
and $\bar{N}_y(n) \ge N_y(n)$ sample path wise. 
Further, we have 
\EQ{
\sum_{j=1}^{\bar{N}_y(n) + 1}\bar{H}_j
= \bar{\tau}_{\bar{N}_y(n)+1} 
= \bar{\tau}_{\bar{N}_y(n)} + \bar{H}_{\bar{N}_y(n)+1}
\le n + M.
}
Applying Wald's Lemma to stopping time $N_y(n)+1$ with respect to \iid sequence $H$ and stopping time $\bar{N}_y(n)+1$ with respect to \iid sequence $\bar{H}$, 
and monotonicity of expectation, we get 
\EQ{
\E_y(N_y(n) + 1)\bar{\mu}_{yy} \le \E_y(\bar{N}_y(n) + 1)\bar{\mu}_{yy} \le n + M. 
}
Taking limits, we obtain $\lim\sup_{n \in \N}\frac{\sum_{k=1}^np_{yy}^{(k)}}{n} \le \frac{1}{\bar{\mu}_{yy}}$. 
Letting $M$ grow arbitrarily large, we obtain the upper bound. 
%From the strong law of large numbers, we know that the following almost sure equality holds, 
%\EQ{
%\lim_{m \in \N}\frac{\sum_{\ell=1}^{m}H_1{m} = \mu_{yy}. 
%}
%We recognise that the number of visits $m$ to state $y$ in $n$ steps is $\sum_{k=1}^n1\set{X_k = y}$, 
%and the sum of hitting times is $\sum_{\ell = 1}^{m}H_1 \le n < \sum_{\ell = 1}^{m+1}H_1$. 
%That is, we have 
%\EQ{
%\frac{m}{\sum_{\ell = 1}^{m+1}H_1 < \frac{\sum_{k=1}^n1\set{X_k = y}}{n} \le \frac{m}{\sum_{\ell = 1}^{m}H_1. 
%}
%We can write the sum $\sum_{k=1}^nP_{xy}^{(k)} = \E_x\sum_{k = 1}^n1\set{X_k = y} = \E_xN_y(n)$. 
%By weak law of large numbers, we get $\lim_{n \to \infty}\frac{\E_xN_y(n)}{n} =  $ 

\item[Starting from $x$: ] Further, we observe that $p_{xy}^{(k)} = \sum_{s=0}^{k-1}f_{xy}^{(k-s)}p_{yy}^{(s)}$.  
Since $1 = f_{xy} = \sum_{k \in \N}f_{xy}^{(k)}$, we have 
\EQ{
\sum_{k=1}^{n}p_{xy}^{(k)} = \sum_{k=1}^n\sum_{s=0}^{k-1}f_{xy}^{(k-s)}p_{yy}^{(s)} = \sum_{s = 0}^{n-1}p_{yy}^{(s)}\sum_{k - s = 1}^{n-s}f_{xy}^{(k-s)}  = \sum_{s = 0}^{n-1}p_{yy}^{(s)} -   \sum_{s = 0}^{n-1}p_{yy}^{(s)}\sum_{k > n-s}f_{xy}^{(k)}. 
%\le \sum_{s = 0}^{n-1}p_{yy}^{(s)}f_{xy}. 
}
Since the series $\sum_{k \in \N}f_{xy}^{(k)}$ converges, we get 
%\EQ{
$\lim_{n \to \infty}\frac{\sum_{k=1}^np_{xy}^{(k)}}{n} = \lim_{n \to \infty}\frac{\sum_{k=1}^np_{yy}^{(k)}}{n}. $
%}
\end{itemize}
\end{proof}
\end{document}
